\documentclass[12pt]{article}
\usepackage[T1]{fontenc}
\usepackage{lmodern}

\usepackage[margin=1in]{geometry}
\usepackage{amsmath,amssymb,amsthm,mathtools}
\usepackage{booktabs,array,enumitem,microtype}
\usepackage[colorlinks=true,allcolors=blue]{hyperref}

\newtheorem{theorem}{Theorem}[section]
\newtheorem{proposition}[theorem]{Proposition}
\newtheorem{corollary}[theorem]{Corollary}
\newtheorem{lemma}[theorem]{Lemma}
\theoremstyle{definition}
\newtheorem{definition}[theorem]{Definition}
\newtheorem{example}[theorem]{Example}
\theoremstyle{remark}
\newtheorem{remark}[theorem]{Remark}

\newcommand{\ii}{\mathrm i}
\newcommand{\RR}{\mathbb R}
\newcommand{\CC}{\mathbb C}
\newcommand{\Herm}{\operatorname{Herm}}
\newcommand{\Pos}{\mathcal P}
\newcommand{\SO}{\operatorname{SO}}
\newcommand{\SU}{\operatorname{SU}}
\newcommand{\tr}{\operatorname{tr}}
\newcommand{\rank}{\operatorname{rank}}
\newcommand{\ip}[2]{\left\langle #1,#2\right\rangle}

\title{Green Geometry of Positive Hermitian Instruments:\\
Hyperbolic Transport, Relational Holonomy, and\\
Fixed Five-Functional Tomography}
\author{Luqman Mahmood\thanks{Copyright 2026 Luqman Mahmood. This paper is licensed under Creative Commons Attribution 4.0 International (CC BY 4.0); see \url{https://creativecommons.org/licenses/by/4.0/}. Software is excluded from this grant.}}
\date{1 October 2026}
\hypersetup{pdftitle={Green Geometry of Positive Hermitian Instruments: Hyperbolic Transport, Relational Holonomy, and Fixed Five-Functional Tomography},pdfauthor={Luqman Mahmood}}

\begin{document}
\maketitle

\begin{abstract}
We study the cone $\Pos$ of positive Hermitian $2\times2$ matrices with the
Green/Riesz metric
$g_G(U,V)=\operatorname{Re}\tr(G^{-1}UG^{-1}V)$.  A direct calculation gives
its Levi--Civita connection and splits the geometry into a flat scale
direction and determinant-one shape leaves of sectional curvature $-1/2$.
Parallel transport around a closed instrument loop acts on the traceless
Hermitian directions by a rotation in $\SO(3)$.  This return is relational: it
is visible against a fixed reference frame and disappears when the reference
is co-transported.  We then ask how few real-linear readings determine the
return rotation.  Three antisymmetric readings give local coordinates, while
adding the trace still fails at half-turns.  Five fixed functionals suffice
globally.  The proof combines six rank-two projections in $\RR^4$, exact
whitening, the quaternion--rotation dictionary, and an exact real-kernel
exclusion.  By the published identity $m_{\RR}(4)=6$, five is also
minimal among fixed real-linear schemes.  For the resulting positive-projector
design we prove a kernel floor $1/1024$, a two-sided rotation-angle estimate,
and condition number at most $560.681$.  A separate rational four-plane has
proved floor $1/32$ and an explicit abstract isometric realization with a
larger lower angle constant.  These are comparisons of proved bounds, not of
unknown true optima; no stochastic-noise, device, or physical-realization conclusion is
drawn.
\end{abstract}

\section{Introduction}\label{sec:introduction}

Positive-definite matrix geometry is classical.  In particular, the
affine-invariant trace metric, its congruence invariance, geodesics, connection,
and nonpositive curvature belong to the established geometry of positive
matrices
\cite{PennecFillardAyache2006,Moakher2005,Bhatia2007,
ThanwerdasPennec2023,HiaiPetz2009,BhatiaHolbrook2006}.  Its covariance
Fisher--Rao geometry is likewise established
\cite{Skovgaard1984,AmariNagaoka2000,LengletEtAl2006}.  Quantum-state
monotone metrics are a distinct comparison family
\cite{PetzSudar1996}; no quantum statistical interpretation is inferred here.
The determinant-one
positive Hermitian $2\times2$ model is the standard hyperbolic
$\mathrm{SL}(2,\CC)/\SU(2)$ model \cite{Bryant1987}.  Likewise, attitude
estimation, rotation synchronization, continuous Euclidean representations
of $\SO(3)$, and real generalized phase retrieval are established subjects
\cite{Wahba1965,WangSinger2013,ZhouEtAl2019,Xu2018,WangXu2019}.

The return matrix $C_{ab}=\tfrac12\tr(\sigma_aH\sigma_bH^\dagger)$ is
also the Pauli transfer matrix of a unitary qubit channel.  The closest
tomographic comparison is therefore process tomography under a unitarity
promise, rather than state tomography
\cite{GutoskiJohnston2014,BaldwinKalevDeutsch2014,DiColandreaEtAl2023}.
Our count concerns exact synthesized real-linear coordinates on $\SO(3)$;
it is not a count of experimental settings, samples, or interactive
observables in those works.

Our question has two parts.  What does an observer who retains a fixed Green
reference frame see after parallel transport around a loop?  And how few
real-linear readings determine that return?  We derive the geometry from the
Green/Riesz pairing, express the return as a relational rotation, and then
solve the associated fixed-functional reconstruction problem.

The mathematical chain is
\[
\begin{array}{ccccc}
\boxed{\text{Green pairing}}&\longrightarrow&
\boxed{\text{instrument metric}}&\longrightarrow&
\boxed{\text{parallel transport}}\\[-1mm]
&&&&\big\downarrow\\[-1mm]
&&\boxed{\text{five fixed functionals}}&\longleftarrow&
\boxed{\text{relational }\SO(3)\text{ return}}
\end{array}.
\]
Every arrow is finite-dimensional and exact.  The geometry is mathematical:
we do not infer an evolution law or physical realization of arbitrary loops.

The paper has five main steps.  First,
Theorem~\ref{thm:green-geometry} derives the metric connection, curvature, and
scale--shape product.  Second, Proposition~\ref{prop:exact-transport} and
Theorem~\ref{thm:relational-holonomy} derive exact transport and relational
holonomy.  Third, Theorem~\ref{thm:local-three} and
Proposition~\ref{prop:half-turn-wall} distinguish local coordinates from a
global wall.  Fourth, Theorems~\ref{thm:kernel-certificate} and
\ref{thm:fixed-five} give a fixed, nonadaptive, globally injective
five-functional construction.  Corollary~\ref{cor:fixed-five-optimality-external}
uses the known phase-retrieval count to obtain fixed-linear optimality, and
Proposition~\ref{prop:adaptive-five} records the adaptive comparison.
 Theorem~\ref{thm:projector-conditioning} and
 Corollary~\ref{cor:projector-conditioning} give quantitative bounds
for the positive-projector construction.  Finally,
Theorem~\ref{thm:kstar-floor} and Corollary~\ref{cor:kstar-conditioning}
record a separately specified rational four-plane with a larger proved
conditioning floor.  The two designs serve different purposes: the first is
tied to positive projections, whereas the second is an abstract measurement
kernel.

\section{The Green metric and hyperbolic shape geometry}
\label{sec:green_metric}

\begin{definition}[Instrument and tangent carrier]\label{def:instrument-space}
Let
\[
 \Pos=\{G\in\Herm(2):G>0\}.
\]
The canonical tangent identification is $T_G\Pos=\Herm(2)$.  Under an
invertible passive complex frame $W$, instrument and tangent transform by the
same congruence,
\begin{equation}\label{eq:passive-congruence}
 G'=W^{-\dagger}GW^{-1},\qquad
 U'=W^{-\dagger}UW^{-1}.
\end{equation}
\end{definition}

\begin{definition}[Green source and reading]\label{def:green-reading}
For $G\in\Pos$ and $Q\in\Herm(2)$, the Green/Riesz source is
$S_Q=G^{-1}Q$.  The scalar reading of a tangent response $U$ is
\begin{equation}\label{eq:green-reading}
 \mathcal E_G(Q,U)
 =\operatorname{Re}\tr(G^{-1}QG^{-1}U).
\end{equation}
We apply the same pairing to tangent variations and write
\begin{equation}\label{eq:green-metric}
 g_G(U,V)=\operatorname{Re}\tr(G^{-1}UG^{-1}V).
\end{equation}
\end{definition}

\begin{lemma}[Determinant tangent and scale split]\label{lem:det-tangent}
For $G\in\Pos$,
\begin{equation}\label{eq:det-derivative}
 d(\det)_G[U]=\det(G)\tr(G^{-1}U).
\end{equation}
Consequently the tangent space to a fixed-determinant leaf is
\[
 T_G\Pos_{\det G}=\{U\in\Herm(2):\tr(G^{-1}U)=0\},
\]
and the exact projection onto it is
\begin{equation}\label{eq:shape-projector}
 \Pi_G(U)=U-\tfrac12\tr(G^{-1}U)G.
\end{equation}
The scale vector $G$ is orthogonal to this leaf and $g_G(G,G)=2$.
\end{lemma}

\begin{proof}
Jacobi's formula gives~\eqref{eq:det-derivative}.  Since
$\tr(G^{-1}G)=2$, equation~\eqref{eq:shape-projector} is trace-free in the
required sense.  Finally,
$g_G(G,U)=\tr(G^{-1}U)$ and $g_G(G,G)=\tr I_2=2$.
\end{proof}

The next theorem derives the normalization of the classical affine-invariant
positive-matrix geometry induced by this Green/Riesz pairing
\cite{PennecFillardAyache2006,Moakher2005,Bhatia2007,
ThanwerdasPennec2023,HiaiPetz2009}.

\begin{theorem}[Green geometry]\label{thm:green-geometry}
The pairing~\eqref{eq:green-metric} is a positive Riemannian metric,
invariant under~\eqref{eq:passive-congruence}.  Its unique torsion-free,
metric-compatible connection is
\begin{equation}\label{eq:levi-civita}
 \nabla_XY
 =DY[X]-\tfrac12\bigl(XG^{-1}Y+YG^{-1}X\bigr).
\end{equation}
With the curvature convention
$R(X,Y)Z=\nabla_X\nabla_YZ-\nabla_Y\nabla_XZ-\nabla_{[X,Y]}Z$,
\begin{align}
R_G(X,Y)Z=-\tfrac14\bigl(
 &XG^{-1}YG^{-1}Z-YG^{-1}XG^{-1}Z \notag\\
 &-ZG^{-1}XG^{-1}Y+ZG^{-1}YG^{-1}X
\bigr).\label{eq:curvature}
\end{align}
The fixed-determinant leaves are totally geodesic.  Every nondegenerate
two-plane tangent to such a leaf has sectional curvature
\begin{equation}\label{eq:shape-curvature}
 K=-\tfrac12.
\end{equation}
If
\begin{equation}\label{eq:scale-shape-coordinate}
 G=e^sP,\qquad \det P=1,
\end{equation}
then the metric splits exactly as
\begin{equation}\label{eq:metric-product}
 g_G=2\,ds^2+g_P.
\end{equation}
Thus the full cone is nonpositively curved but not constant-curvature:
scale--shape planes are flat.
\end{theorem}

\begin{proof}
Put $A=G^{-1}$ and
$H=A^{1/2}UA^{1/2}$.  Then $H$ is Hermitian and
\[
 g_G(U,U)=\tr(H^2)>0
\]
for $U\ne0$.  Under~\eqref{eq:passive-congruence},
$(G')^{-1}=WG^{-1}W^\dagger$; cyclicity of trace proves invariance.

To derive, rather than import, the connection, use
$D_XA=-AXA$.  For constant Hermitian vector fields,
\[
D_Xg(Y,Z)
=-\operatorname{Re}\tr(AXAYAZ+AYAXAZ).
\]
Insert this identity and its cyclic permutations into the Koszul formula.
For the general symmetric one-coefficient bilinear ansatz
$\Gamma_G(X,Y)=\alpha(XAY+YAX)$, cyclic trace reduction leaves
$2\alpha+1=0$.  Hence $\alpha=-1/2$, which gives
~\eqref{eq:levi-civita}.  Symmetry in $X,Y$ gives zero torsion; substituting
$D_XA=-AXA$ gives $\nabla g=0$.  Uniqueness is the standard uniqueness part
of the Koszul construction.

Differentiate~\eqref{eq:levi-civita}, again using $D_XA=-AXA$, and subtract
the expression with $X,Y$ exchanged.  The derivative terms cancel and leave
~\eqref{eq:curvature}.  Whitening by $G^{-1/2}$ reduces it to
\begin{equation}\label{eq:whitened-curvature}
 R(X,Y)Z=-\tfrac14[[X,Y],Z],
\end{equation}
and therefore
\begin{equation}\label{eq:sectional-numerator}
 \ip{R(X,Y)Y}{X}
 =\tfrac14\tr([X,Y]^2)\le0.
\end{equation}
For traceless Hermitian $2\times2$ matrices,
\[
 -\tr([X,Y]^2)
 =2\{\tr(X^2)\tr(Y^2)-\tr(XY)^2\}.
\]
Division by the Gram determinant in~\eqref{eq:sectional-numerator} gives
$K=-1/2$.

If tangent fields $X,Y$ obey $\tr(AX)=\tr(AY)=0$, differentiating the first
constraint in the $Y$ direction and using~\eqref{eq:levi-civita} shows
$\tr(A\nabla_YX)=0$.  Hence the leaf is totally geodesic.  Finally,
$G^{-1}dG=ds\,I+P^{-1}dP$ in the trace calculation, and
$\tr(P^{-1}dP)=d\log\det P=0$.  The cross term vanishes and $\tr I_2=2$,
giving~\eqref{eq:metric-product}.
\end{proof}

\begin{remark}[Classical identification]\label{rem:classical-geometry}
Theorem~\ref{thm:green-geometry} recovers the classical affine-invariant
positive-matrix metric \cite{ThanwerdasPennec2023,HiaiPetz2009}.  On
$\det P=1$, the conventional hyperbolic metric has curvature $-1$; the
normalization~\eqref{eq:green-metric} is twice that metric and therefore has
curvature $-1/2$ \cite{Bryant1987}.  No new affine-invariant metric or
hyperbolic model is claimed.  Geometrically, radial motion changes only the
overall scale, while motion inside a determinant-one leaf follows the
three-dimensional hyperbolic shape geometry.  Curvature then predicts the
leading small-loop return: an infinitesimal loop of oriented area $A$ in an
orthonormal shape plane produces a rotation of order $A/2$.
\end{remark}

\section{Exact geodesics and parallel transport}\label{sec:transport}

Affine-invariant positive-matrix geodesics are classical
\cite{BhatiaHolbrook2006}; this section derives the endpoint formula used by
the Green normalization above.

The geodesic equation obtained by setting
$\nabla_{\dot G}\dot G=0$ in~\eqref{eq:levi-civita} is
\begin{equation}\label{eq:geodesic-equation}
 \ddot G=\dot G\,G^{-1}\dot G.
\end{equation}
The parallel-field equation along the same path is
\begin{equation}\label{eq:parallel-ode}
 \dot U=\tfrac12\bigl(\dot G\,G^{-1}U+UG^{-1}\dot G\bigr).
\end{equation}

\begin{proposition}[Exact endpoint transport]\label{prop:exact-transport}
Let $G,H\in\Pos$ satisfy $\det G=\det H$.  Put
\[
 M=HG^{-1},\qquad
 E_{G\to H}=\frac{M+I}{\sqrt{\tr(M)+2}},
\]
where the square root in the denominator is the positive one.  Then
\begin{equation}\label{eq:transport-identities}
 E_{G\to H}^{\,2}=M,\quad
 E_{G\to H}G=GE_{G\to H}^{\dagger},\quad
 E_{G\to H}GE_{G\to H}^{\dagger}=H,\quad
 \det E_{G\to H}=1.
\end{equation}
Parallel transport along the affine-invariant geodesic from $G$ to $H$ is
\begin{equation}\label{eq:parallel-transport}
 \mathcal T_{G\to H}(U)=E_{G\to H}UE_{G\to H}^{\dagger}.
\end{equation}
It preserves Hermiticity, the determinant-tangent constraint, and the metric,
and its reverse is $\mathcal T_{H\to G}$.  Segment transports compose in
chronological order.  Under the passive congruence
\eqref{eq:passive-congruence}, the transport is covariant:
\begin{equation}\label{eq:transport-covariance}
 \mathcal T_{G'\to H'}(U')
 =W^{-\dagger}\mathcal T_{G\to H}(U)W^{-1}.
\end{equation}
Equal determinants are needed for the displayed closed form and for the
path to stay on one determinant leaf.  For arbitrary positive $G,H$, the
same transport formula holds with $E_{G\to H}=(HG^{-1})^{1/2}$, using the
positive spectral square root.
\end{proposition}

\begin{proof}
$M$ is similar to the positive Hermitian matrix
$G^{-1/2}HG^{-1/2}$, so its eigenvalues are positive.  Equal determinants
give $\det M=1$.  Cayley--Hamilton then yields
$M^2-(\tr M)M+I=0$, hence
$(M+I)^2=(\tr M+2)M$ and the first identity in
~\eqref{eq:transport-identities}.  Since $MG=H=GM^\dagger$, every real
polynomial in $M$ satisfies the analogous $G$-adjoint identity; this proves
the second identity.  The third follows from
$EGE^\dagger=E^2G=MG=H$, and the positive square root has determinant one.

Define $G(t)=M^tG$ by the positive spectral calculus.  Direct differentiation
gives~\eqref{eq:geodesic-equation}.  Substitution shows that
$U(t)=M^{t/2}U(0)(M^\dagger)^{t/2}$ solves~\eqref{eq:parallel-ode};
at $t=1$ it is~\eqref{eq:parallel-transport}.  The remaining assertions
follow by congruence, cyclicity of trace, composition of the segment maps,
and reversal of the geodesic.  In particular,
$M'=W^{-\dagger}MW^\dagger$ and
$E_{G'\to H'}=W^{-\dagger}E_{G\to H}W^\dagger$, which gives
\eqref{eq:transport-covariance}.
\end{proof}

\section{Closed loops and relational Green holonomy}\label{sec:holonomy}

Holonomy of symmetric spaces is classical
\cite{KobayashiNomizu1963,AmbroseSinger1953}.  Berry-phase anholonomy is a
useful neighboring comparison \cite{Berry1984}, but the return below is a
relational Green-frame rotation; it is not a physical quantum phase or an
adiabatic-time claim.

\begin{definition}[Relational return]\label{def:relational-return}
Let $G_0,\ldots,G_n$ be positive determinant-one instruments with
$G_n=G_0=I$.  In chronological order put
\begin{equation}\label{eq:loop-spin}
 H_{\mathcal L}
 =E_{G_{n-1}\to G_n}\cdots E_{G_1\to G_2}E_{G_0\to G_1}.
\end{equation}
For the Pauli frame $(\sigma_1,\sigma_2,\sigma_3)$, define
\begin{equation}\label{eq:return-action}
 C_{ab}(\mathcal L)
 =\tfrac12\tr\!\left(
 \sigma_aH_{\mathcal L}\sigma_bH_{\mathcal L}^{\dagger}\right).
\end{equation}
\end{definition}

Because~\eqref{eq:loop-spin} returns the base instrument,
$H_{\mathcal L}H_{\mathcal L}^{\dagger}=I$ and
$\det H_{\mathcal L}=1$.  Thus $H_{\mathcal L}\in\SU(2)$ and
$C(\mathcal L)\in\SO(3)$.  The rotation-matrix readout lies within the
classical literature on attitude estimation and rotation synchronization
\cite{Wahba1965,WangSinger2013}; the statement here concerns its exact
derivation from retained versus transported Green reference frames.

\begin{theorem}[Relational holonomy and source readout]
\label{thm:relational-holonomy}
For every determinant-one loop based at $I$,
\begin{equation}\label{eq:fixed-green-readout}
 C_{ab}(\mathcal L)
 =\tfrac12 g_I\!\left(
 \sigma_a,H_{\mathcal L}\sigma_bH_{\mathcal L}^{\dagger}\right).
\end{equation}
If both the query and reference frames are co-transported, their mutual
reading is instead
\begin{equation}\label{eq:cotransported-readout}
 \tfrac12g_I(H_{\mathcal L}\sigma_aH_{\mathcal L}^{\dagger},
             H_{\mathcal L}\sigma_bH_{\mathcal L}^{\dagger})
 =\delta_{ab}.
\end{equation}
Loop reversal gives $C^{-1}$, concatenation gives the ordered product, and
$H_{\mathcal L}$ and $-H_{\mathcal L}$ have the same readout.  Noncollinear
loops can have $C\ne I$, while collinear geodesic loops return $I$.  The
infinitesimal return algebra on the three Pauli modalities is all of
$\mathfrak{so}(3)$.
\end{theorem}

\begin{proof}
At $G=I$, equation~\eqref{eq:green-metric} is the Hilbert--Schmidt pairing,
so~\eqref{eq:fixed-green-readout} is exactly
~\eqref{eq:return-action}.  Metric preservation proves
~\eqref{eq:cotransported-readout}.  Functoriality of parallel transport gives
the reversal and concatenation laws, and the adjoint action is blind to the
center of $\SU(2)$.

For traceless Hermitian $X,Y$, the small-loop return generator is
$[X,Y]/4$ with the boundary orientation fixed above.  Its adjoint action on
the Pauli frame is a real skew $3\times3$ matrix.  The three independent
Pauli pairs produce three independent generators and close under commutator
as $\mathfrak{so}(3)$.  This is the local holonomy algebra, consistent with
the classical curvature--holonomy relation \cite{AmbroseSinger1953}.
Exact finite nonidentity and collinear examples are displayed below.
\end{proof}

Here are three exact loops realizing nontrivial returns.  For $a,b>0$ and
$|z|=1$, write
\[
 \mathcal D(a)=\begin{pmatrix}a^2&0\\0&a^{-2}\end{pmatrix},\qquad
 \mathcal B(b,z)=\begin{pmatrix}c_b&z s_b\\\bar z s_b&c_b\end{pmatrix},
 \quad
 c_b=\frac{b^2+b^{-2}}2,\quad s_b=\frac{b^2-b^{-2}}2.
\]
Both matrices are positive Hermitian with determinant one.

\begin{example}[A real-axis triangle]\label{ex:holonomy-real}
The geodesic triangle
$I\to\mathcal D(7/3)\to\mathcal B(5/3,1)\to I$ has return
\begin{equation}\label{eq:return-real-axis}
C_1=
\begin{pmatrix}
217449/268649&0&-157760/268649\\
0&1&0\\
157760/268649&0&217449/268649
\end{pmatrix}.
\end{equation}
It fixes the second Pauli axis.
\end{example}

\begin{example}[An imaginary-axis triangle]\label{ex:holonomy-imaginary}
The triangle
$I\to\mathcal D(8/5)\to\mathcal B(9/4,\ii)\to I$ has return
\begin{equation}\label{eq:return-imaginary-axis}
C_2=
\begin{pmatrix}
1&0&0\\
0&34051232/40477457&21884655/40477457\\
0&-21884655/40477457&34051232/40477457
\end{pmatrix}.
\end{equation}
It fixes the first Pauli axis.
\end{example}

\begin{example}[A generic complex-axis triangle]\label{ex:holonomy-generic}
For $z=(3+4\ii)/5$, the triangle
$I\to\mathcal D(11/7)\to\mathcal B(6/5,z)\to I$ has return
\begin{equation}\label{eq:return-generic-axis}
C_3=
\begin{pmatrix}
673203337/676026025&3763584/676026025&-2463912/27041041\\
3763584/676026025&671007913/676026025&3285216/27041041\\
2463912/27041041&-3285216/27041041&26727409/27041041
\end{pmatrix}.
\end{equation}
Its fixed axis is not a coordinate Pauli axis.  By contrast, every collinear
loop $I\to\mathcal D(a)\to\mathcal D(b)\to I$ has identity return.
\end{example}

\section{Local tomography and the half-turn wall}
\label{sec:local_tomography}

Minimal tomography for trace-one qubits and informational completeness under
prior information are established, with measurement objects and counts that
differ from the unnormalized positive-cone Green-response model used here
\cite{RehacekEtAl2004,HeinosaariMazzarellaWolf2013}.  Consequently, every
minimality statement below is restricted to its declared reading scheme or
to the separately stated class of fixed real-linear functionals; it is not a
general tomography lower bound.

\begin{definition}[Antisymmetric local reading]\label{def:local-reading}
For $C\in\SO(3)$, set
\begin{equation}\label{eq:w-reading}
 w(C)=\tfrac12
 \begin{pmatrix}
 C_{32}-C_{23}\\ C_{13}-C_{31}\\ C_{21}-C_{12}
 \end{pmatrix}.
\end{equation}
For $w\in\RR^3$, let $[w]_\times x=w\times x$.
\end{definition}

\begin{theorem}[Three local readings]\label{thm:local-three}
The differential of~\eqref{eq:w-reading} at the identity has rank three.
On the branch
\[
 c=\tfrac12(\tr C-1)>0,
\]
the three readings determine $C$ exactly by
\begin{equation}\label{eq:local-rodrigues}
 c=+\sqrt{1-w^\mathsf Tw},\qquad
 C=cI+[w]_\times+\frac{ww^\mathsf T}{1+c}.
\end{equation}
No continuous scalar instrument with fewer than three outputs is locally
injective on an open subset of the three-dimensional return manifold.
\end{theorem}

\begin{proof}
Write a rotation near the identity as
$C=\exp([\omega]_\times)$.  Then $w(C)=\omega+O(|\omega|^2)$, proving the
rank statement.  Rodrigues' formula gives
$w=\sin\theta\,n$ and $c=\cos\theta$; on $c>0$ the positive square-root
choice is forced and rearrangement yields~\eqref{eq:local-rodrigues}.
Exact multiplication verifies $C^\mathsf TC=I$, $\det C=1$, and readback.
The lower bound is invariance of dimension: an open subset of a
three-manifold cannot embed continuously into $\RR^m$ for $m<3$.
\end{proof}

\begin{proposition}[Natural four-reading wall]\label{prop:half-turn-wall}
The four readings $(w(C),\tr C)$ are not globally injective.  Every half-turn
about a unit axis $n$ has
\begin{equation}\label{eq:half-turn}
 C_n=2nn^\mathsf T-I,\qquad
 (w(C_n),\tr C_n)=(0,0,0,-1).
\end{equation}
\end{proposition}

\begin{proof}
$C_n$ is symmetric, so its antisymmetric reading is zero.  Its eigenvalues
are $(1,-1,-1)$, so its trace is $-1$.  Distinct axes not related by sign
give distinct rotations with the same four-reading code.
\end{proof}

\section{Five fixed real-linear Green functionals}
\label{sec:fixed_five}

\begin{definition}[Fixed functional]\label{def:fixed-functional}
A fixed real-linear Green functional is
\[
 \ell_K(C)=\ip{K}{C}_F=\tr(K^\mathsf TC),
 \qquad K\in\RR^{3\times3},
\]
where $K$ is selected before $C$ is known.  Because the nine entries of
$C$ are the fixed-reference readings~\eqref{eq:fixed-green-readout},
$\ell_K$ lies in their real-linear span.  Fixed excludes data-dependent
pivots or a change of $K$ after observing $C$.
\end{definition}

Use quaternion order $q=(w,x,y,z)^\mathsf T$ and the standard double cover
$q\mapsto R(q)$:
\begin{equation}\label{eq:quaternion-rotation}
R(q)=
\begin{pmatrix}
w^2+x^2-y^2-z^2&2(xy-wz)&2(xz+wy)\\
2(xy+wz)&w^2-x^2+y^2-z^2&2(yz-wx)\\
2(xz-wy)&2(yz+wx)&w^2-x^2-y^2+z^2
\end{pmatrix}.
\end{equation}
For unit $q$, $R(q)\in\SO(3)$ and $R(q)=R(-q)$.

\begin{lemma}[Quaternion--rotation dictionary]
\label{lem:quaternion-dictionary}
Let $D=(d_{\mu\nu})_{\mu,\nu=0}^3$ be real symmetric and trace-free.  Define
$K(D)\in\RR^{3\times3}$ by
\begin{align}
K_{11}&=(d_{00}+d_{11})/2,&
K_{22}&=(d_{00}+d_{22})/2,&
K_{33}&=(d_{00}+d_{33})/2,\notag\\
K_{23}&=(d_{23}-d_{01})/2,&
K_{32}&=(d_{23}+d_{01})/2,\notag\\
K_{13}&=(d_{02}+d_{13})/2,&
K_{31}&=(d_{13}-d_{02})/2,\label{eq:inverse-dictionary}\\
K_{12}&=(d_{12}-d_{03})/2,&
K_{21}&=(d_{12}+d_{03})/2.\notag
\end{align}
Then, coefficientwise for every $q\in\RR^4$,
\begin{equation}\label{eq:dictionary-readback}
 q^\mathsf TDq=\ip{K(D)}{R(q)}_F.
\end{equation}
Together with $q^\mathsf Tq$, the nine quadratic entries of $R(q)$ span
$\operatorname{Sym}(4)$ with rank ten.
\end{lemma}

\begin{proof}
Substitute~\eqref{eq:quaternion-rotation} into
~\eqref{eq:dictionary-readback}.  Direct coefficient comparison gives
$d_{00},d_{11},d_{22},d_{33}$ on the diagonal, using $\tr D=0$, and
$2d_{\mu\nu}$ on each cross term.  This proves the identity.  The inverse recovers
all nine trace-free coordinates, while the norm recovers the scalar
direction, giving rank ten.
\end{proof}

The fixed construction starts from Xu's six rank-two projections
\cite[Theorem~4.2]{Xu2018}, represented here by the columns $u_j,v_j$ of
\begin{equation}\label{eq:exact-uv}
U=\begin{pmatrix}
1&0&0&0&1&1\\
0&1&0&0&0&1\\
0&0&1&0&-4&2\\
0&0&0&1&-3&3
\end{pmatrix},\qquad
V=\begin{pmatrix}
0&5&-1&-1&-17&-5\\
-5&0&0&5&4&4\\
2&-2&0&-3&-2&2\\
-2&1&0&0&-3&-1
\end{pmatrix}.
\end{equation}
Define
\begin{equation}\label{eq:projectors}
 P_j=\frac{u_ju_j^\mathsf T}{u_j^\mathsf Tu_j}
     +\frac{v_jv_j^\mathsf T}{v_j^\mathsf Tv_j},
\qquad S=\sum_{j=1}^6P_j.
\end{equation}
The exact factorization used below is $S=LDL^\mathsf T$, with
\begin{equation}\label{eq:L-factor}
L=\begin{pmatrix}
1&0&0&0\\
-1206036/7353455&1&0&0\\
-629974/7353455&-620743491/3418884916&1&0\\
865837/7353455&562593553/3418884916&
39888940419/121436249827&1
\end{pmatrix}
\end{equation}
and
\begin{equation}\label{eq:D-factor}
D=\operatorname{diag}\!\left(
\frac{1470691}{332787},
\frac{3418884916}{1213320075},
\frac{121436249827}{51283273740},
\frac{213445555084}{121436249827}
\right).
\end{equation}
Put
\begin{equation}\label{eq:whitening}
 C_{\mathrm w}=D^{-1/2}L^{-1},\qquad
 B_j=C_{\mathrm w}P_jC_{\mathrm w}^{\mathsf T}.
\end{equation}

\begin{lemma}[Exact projector and whitening identities]\label{lem:projectors}
\label{lem:whitening}
For every $j$, $u_j^\mathsf Tv_j=0$, $P_j=P_j^\mathsf T=P_j^2$,
$P_j\ge0$, and $\rank P_j=2$.  The four diagonal entries of~\eqref{eq:D-factor}
are positive, $S=LDL^\mathsf T$, and
\begin{equation}\label{eq:tight-sum}
 \sum_{j=1}^6B_j=I_4.
\end{equation}
\end{lemma}

\begin{proof}
All claims before the factorization follow directly from the six integer
column pairs in~\eqref{eq:exact-uv}; orthogonality makes
~\eqref{eq:projectors} the sum of two orthogonal rank-one projections.
Every diagonal entry of $D$ has positive numerator and denominator.
Rational multiplication verifies $S=LDL^\mathsf T$.  Therefore
\[
\sum_jB_j
=D^{-1/2}L^{-1}(LDL^\mathsf T)L^{-\mathsf T}D^{-1/2}=I_4.
\]
\end{proof}

For each $j$, write
\begin{equation}\label{eq:Dj-alpha}
 \alpha_j=\tfrac14\tr B_j,\qquad
 D_j=B_j-\alpha_jI_4,\qquad K_j=K(D_j).
\end{equation}
The five promised functionals are
\begin{equation}\label{eq:functionals}
 \ell_j(C)=\ip{K_j}{C}_F,\qquad j=1,\ldots,5.
\end{equation}
Equations~\eqref{eq:exact-uv}--\eqref{eq:Dj-alpha} are the exact definition of
every coefficient; no decimal or fitted parameter enters.

\section{Exact rank-two-kernel reconstruction and global injectivity}
\label{sec:certificate}

\begin{lemma}[Real phase-retrieval kernel criterion]
\label{lem:rank-two-kernel}
Let $B_1,\ldots,B_m$ be real symmetric matrices with
$\sum_jB_j=I$.  The quadratic map
\[
 q\longmapsto(q^\mathsf TB_1q,\ldots,q^\mathsf TB_mq)
\]
determines every $q\in\RR^4$ up to sign if and only if its symmetric kernel
contains no nonzero matrix of rank at most two.
\end{lemma}

\begin{proof}
If $q$ and $r$ have equal measurements, then
$Q=qq^\mathsf T-rr^\mathsf T$ is a symmetric kernel element of rank at most
two.  Conversely, let a nonzero symmetric kernel element $Q$ have rank at
most two.  Since $\sum B_j=I$, $\tr Q=0$.  It cannot be nonzero semidefinite,
so its spectral form is $\lambda aa^\mathsf T-\lambda bb^\mathsf T$ with
$\lambda>0$ and orthonormal $a,b$.  The two non-antipodal vectors
$\sqrt\lambda a,\sqrt\lambda b$ then have equal quadratic measurements.
\end{proof}

\begin{theorem}[Exact six-map kernel exclusion]
\label{thm:kernel-certificate}
For the six projectors~\eqref{eq:projectors}, the linear map
\[
 \mathcal M(Q)=(\tr(P_1Q),\ldots,\tr(P_6Q)),
 \qquad Q\in\operatorname{Sym}(4),
\]
has rank six and a four-dimensional kernel, and that kernel contains no
nonzero real matrix of rank at most two.  Consequently both
$(P_j)_{j=1}^6$ and their congruence-whitened family
$(B_j)_{j=1}^6$ perform real phase retrieval in dimension four.
\end{theorem}

\begin{proof}
Use symmetric coordinate order
\[
(q_{00},q_{01},q_{02},q_{03},q_{11},q_{12},q_{13},q_{22},q_{23},q_{33}).
\]
Exact row reduction of the $6\times10$ measurement matrix gives rank six and
the four rational kernel vectors printed in the exact supplement.  Write a
general kernel matrix as
\[
 Q(a,b,c,d)=aQ_a+bQ_b+cQ_c+dQ_d.
\]
A symmetric $4\times4$ matrix has rank at most two exactly when all ten
upper-triangular entries of its adjugate (the corresponding $3\times3$
cofactors of the matrix) vanish.  Homogeneity permits the
projective cover
\[
d=1,\quad d=0,c=1,\quad d=c=0,b=1,\quad d=c=b=0,a=1.
\]
On $d=1$, exact lexicographic elimination gives a three-element basis and a
degree-ten univariate eliminant in $c$.  Its complete integer coefficients
are printed in the supplement.  A signed-remainder Sturm chain has degree
sequence
\[
10,9,8,7,6,5,4,3,2,1,0
\]
and five sign variations at each of $-\infty$ and $+\infty$, hence no real
root.  The $d=0,c=1$ and $d=c=0,b=1$ ideals both contain $1$.  On the last
patch, clear denominators from each cofactor polynomial, divide its integer
coefficient content, and fix its sign.  This nonzero rational rescaling
preserves the common determinantal zero locus.  The first declared nonzero
primitive denominator-cleared cofactor polynomial then evaluates to
\[
2324385247266000.
\]
Thus the projective cofactor system has no real point, proving the kernel
claim.  Invertible congruence preserves rank and transports kernels between
$P_j$ and $B_j$, so the whitened family has the same property.
\end{proof}

\begin{theorem}[Five fixed-function global tomography]\label{thm:fixed-five}
The fixed map
\begin{equation}\label{eq:fixed-map}
 \Phi:\SO(3)\longrightarrow\RR^5,\qquad
 \Phi(C)=(\ell_1(C),\ldots,\ell_5(C)),
\end{equation}
defined by~\eqref{eq:exact-uv}--\eqref{eq:functionals} is injective.
All five $K_j$ have exact rank three and nonzero determinant.
\end{theorem}

\begin{proof}
Let $C=R(q)$ with $q^\mathsf Tq=1$.  By
Lemma~\ref{lem:quaternion-dictionary},
\begin{equation}\label{eq:y-readings}
 y_j:=q^\mathsf TB_jq=\alpha_j+\ell_j(C).
\end{equation}
The five values in~\eqref{eq:fixed-map} therefore recover
$y_1,\ldots,y_5$.  Tightness~\eqref{eq:tight-sum} and unit norm give
\begin{equation}\label{eq:sixth-recovery}
 y_6=1-\sum_{j=1}^5y_j.
\end{equation}
Theorem~\ref{thm:kernel-certificate} says that the six $y_j$ determine $q$
up to the unavoidable sign $q\sim-q$.  Since $R(q)=R(-q)$, they determine
$C$ uniquely.  Exact determinant evaluation gives the five nonzero values
listed in the supplement, proving the rank statement.
\end{proof}

\begin{remark}[What five counts]\label{rem:not-shots}
The five objects in~\eqref{eq:functionals} are synthesized real-linear
functionals of the nine primitive matrix-entry readings.  Each coefficient
matrix has rank three.  The sum of those ranks, $15$, is only an algebraic
upper bound on a rank-one coefficient decomposition.  It is neither a
primitive Green-shot count nor an optimal implementation cost.
\end{remark}

\subsection{Kernel floors and angle constants}\label{sec:conditioning-definitions}

Let $\mathcal S_0$ be the space of trace-zero real symmetric $4\times4$
matrices with Frobenius inner product, and let
\begin{equation}\label{eq:pair-orbit}
 \mathcal W=\{aa^{\mathsf T}-bb^{\mathsf T}:a,b\in\RR^4,
 \ a^{\mathsf T}a=b^{\mathsf T}b=1,\ a^{\mathsf T}b=0\}.
\end{equation}
Every $Z\in\mathcal W$ has $\|Z\|_F^2=2$.  For a four-plane
$K\subset\mathcal S_0$, with Frobenius projection $\Pi_K$, define its
\emph{kernel floor} by
\begin{equation}\label{eq:kernel-floor-definition}
 m(K)=\min_{Z\in\mathcal W}\|\Pi_{K^\perp}Z\|_F^2
     =2-\max_{Z\in\mathcal W}\|\Pi_KZ\|_F^2.
\end{equation}
This is the normalization used for both four-planes below.  If
$K_1,\ldots,K_4$ is a basis of $K$, put
$\Gamma_{ij}=\tr(K_iK_j)$ and
$t_K(Z)_i=\tr(K_iZ)$.  Then
\begin{equation}\label{eq:projection-height}
 \|\Pi_KZ\|_F^2=t_K(Z)^{\mathsf T}\Gamma^{-1}t_K(Z).
\end{equation}
Equivalently, if $S_K(c)=\sum_i c_iK_i$ and
$\operatorname{spr}(A)=\lambda_{\max}(A)-\lambda_{\min}(A)$, the
spread-pencil identity is
\begin{equation}\label{eq:spread-pencil}
 m(K)=2-\max_{c\ne0}
 \frac{\operatorname{spr}(S_K(c))^2}{c^{\mathsf T}\Gamma c}.
\end{equation}
Indeed, maximizing $\tr(AZ)$ over $Z\in\mathcal W$ gives
$\operatorname{spr}(A)$, and duality with~\eqref{eq:projection-height}
gives~\eqref{eq:spread-pencil}.

For symmetric coordinate matrices $A_1,\ldots,A_5$, write
\begin{equation}\label{eq:realized-objective}
 F_A(a,b)=\sum_{j=1}^5
 \tr\!\bigl(A_j(aa^{\mathsf T}-bb^{\mathsf T})\bigr)^2,
 \qquad (a,b)\text{ orthonormal},
\end{equation}
and let $m_A$ and $M_A$ be its minimum and maximum.
For an injective map $\Psi$, define
\[
c_\Psi=\inf_{R_1\ne R_2}\frac{\|\Psi(R_1)-\Psi(R_2)\|}{d_\angle(R_1,R_2)},
\quad C_\Psi=\sup_{R_1\ne R_2}\frac{\|\Psi(R_1)-\Psi(R_2)\|}{d_\angle(R_1,R_2)},
\quad \kappa_\Psi=C_\Psi/c_\Psi
\]
when $c_\Psi>0$.

\begin{lemma}[Rotation-angle reduction]\label{lem:angle-reduction}
Suppose a map $\Psi:\SO(3)\to\RR^5$ and symmetric matrices
$A_1,\ldots,A_5$ satisfy
\begin{equation}\label{eq:coordinate-difference-identity}
 \Psi(R(q))-\Psi(R(r))
 =\left(\tr\!\bigl(A_j(qq^{\mathsf T}-rr^{\mathsf T})\bigr)\right)_{j=1}^5
\end{equation}
for all unit quaternions $q,r$.  If $d_\angle$ is the rotation angle, then
\begin{equation}\label{eq:generic-angle-reduction}
 \frac{\sqrt{m_A}}{\pi}d_\angle(R_1,R_2)
 \leq\|\Psi(R_1)-\Psi(R_2)\|
 \leq\frac{\sqrt{M_A}}2d_\angle(R_1,R_2).
\end{equation}
\end{lemma}

\begin{proof}
Choose unit quaternions $q,r$ for $R_1,R_2$ and let
$\theta\in[0,\pi/2]$ be their line angle.  A two-dimensional spectral
decomposition gives
 $qq^{\mathsf T}-rr^{\mathsf T}=\sin\theta\,(aa^{\mathsf T}-bb^{\mathsf T})$
 for an orthonormal pair $(a,b)$, while $d_\angle=2\theta$.  Hence
 $\|\Psi(R(q))-\Psi(R(r))\|^2=\sin^2\theta\,F_A(a,b)$, and the bounds follow
 from $2\theta/\pi\leq\sin\theta\leq\theta$.  Conversely, every direction
 $aa^{\mathsf T}-bb^{\mathsf T}$ in~\eqref{eq:pair-orbit} occurs as the
 normalized tangent direction of such a pair difference, so the constants
 $m_A$ and $M_A$ govern the full tangent cone rather than only a subset of it.
\end{proof}

\section{Conditioning of the positive-projector design}
\label{sec:projector-conditioning}

The preceding injectivity construction can now be quantified in the same
finite-dimensional model.  Let $D$ be the positive diagonal factor
in~\eqref{eq:D-factor}.  Let $\widehat K_{\rm proj}$ denote the
four-plane obtained by whitening the symmetric kernel with the $D$-weighted
trace pairing, and let $G_1,\ldots,G_5$ be the exact rational representatives
of the five trace-projected functionals in those coordinates.  Their rational
frame Gram is
\begin{equation}\label{eq:projector-frame-gram}
 (H_{\rm proj})_{jk}=\tr(G_jD^{-1}G_kD^{-1})
 -\frac14\tr(G_jD^{-1})\tr(G_kD^{-1}).
\end{equation}
Write $F_{\rm proj}$ for the realized objective~\eqref{eq:realized-objective}
in these coordinates and set
\begin{align}
 s_{\rm lo}
 &=\frac{812873440394747353670787185015095602549}
 {105051913151663586578335579326020229005312},\label{eq:slo}\\
 L_0=\frac{s_{\rm lo}}{1024}
 &=\frac{812873440394747353670787185015095602549}
 {107573159067303512656215633229844714501439488},\label{eq:L0}\\
 M_0&=\frac{270988332696903}{281474976710656}.\label{eq:M0}
\end{align}

\begin{theorem}[Positive-projector kernel and frame-factor floor]
\label{thm:projector-conditioning}
For the positive-projector construction,
\[
 m(\widehat K_{\rm proj})\geq\frac1{1024},
 \qquad s_{\rm lo}<\lambda_{\min}(H_{\rm proj}),
 \qquad F_{\rm proj}\geq
 \lambda_{\min}(H_{\rm proj})m(\widehat K_{\rm proj})\geq L_0.
\]
In particular, the realized five-functional objective has
$\min F_{\rm proj}\geq L_0\approx7.556470846841732\times10^{-6}$.
\end{theorem}

\begin{proof}
The inequality $m(\widehat K_{\rm proj})\geq1/1024$ is proved by exact
branch-and-bound in rational arithmetic.  The compact coefficient domain in
\eqref{eq:spread-pencil} is covered by finitely many boxes.  Rational
quadratic bounds discard boxes outside the normalized coefficient shell.
Vertex-cone bounds close every remaining box in the supplied cover.  Vertex
eigenvalue bounds use exact integer characteristic polynomials and outward
derivative-test bisection.  The complete replay supplies all $87\,525$
closed and $4\,626$ pruned leaves, with no unresolved box.  The supplementary
proof guide explains the convexity step and the replay checks the entire
subdivision, not merely the earlier one-in-$100$ sample.

The characteristic polynomial of the rational matrix $H_{\rm proj}$ and a
Sturm root count give $H_{\rm proj}>0$ and the strict bracket
$s_{\rm lo}<\lambda_{\min}(H_{\rm proj})$.  The realized five coordinates are
a frame analysis of $\Pi_{\widehat K_{\rm proj}^{\perp}}Z$; the Rayleigh
quotient for~\eqref{eq:projector-frame-gram} therefore gives
\[
 F_{\rm proj}\geq\lambda_{\min}(H_{\rm proj})
 \|\Pi_{\widehat K_{\rm proj}^{\perp}}Z\|_F^2.
\]
Taking the minimum over $\mathcal W$ and substituting the two rational lower
bounds yields~\eqref{eq:L0}.
\end{proof}

\begin{corollary}[Positive-projector rotation-angle stability]
\label{cor:projector-conditioning}
For the map $\Phi$ in~\eqref{eq:fixed-map} and rotation-angle
distance $d_\angle$,
\begin{equation}\label{eq:projector-angle-bounds}
 \frac{\sqrt{L_0}}{\pi}d_\angle(R_1,R_2)
 \leq\|\Phi(R_1)-\Phi(R_2)\|
 \leq\frac{\sqrt{M_0}}2d_\angle(R_1,R_2).
\end{equation}
Consequently
\[
 c_{\rm proj}\geq\frac{\sqrt{L_0}}\pi
 \approx8.750031830575377\times10^{-4},\qquad
 \kappa_{\rm proj}\leq\frac\pi2\sqrt{\frac{M_0}{L_0}}
 \approx560.6809131191.
\]
\end{corollary}

\begin{proof}
A separate Sturm calculation supplies outward eigenvalue brackets for the
positive matrices $B_j$.  Since
$|a^{\mathsf T}B_ja-b^{\mathsf T}B_jb|\leq\lambda_{\max}(B_j)$,
$M_A\leq\sum_{j=1}^5\lambda_{\max}(B_j)^2\leq M_0$.
This is independent of the frame-factor lower bound.
Apply Lemma~\ref{lem:angle-reduction} with
$m_A\geq L_0$.  The result uses one-sided bounds; it does not identify the
unknown true extrema.
\end{proof}

\section{A separate rational conditioning construction}
\label{sec:kstar-conditioning}

It is natural to ask whether a different four-plane improves the floor of
Theorem~\ref{thm:projector-conditioning}.  The supplementary data specify four
rational trace-zero symmetric matrices whose span is a four-plane
$K^\ast\subset\mathcal S_0$.  Unlike $\widehat K_{\rm proj}$, this plane is
not obtained from positive rank-two projections; it is an abstract
measurement design.  Its floor is nevertheless defined by exactly
\eqref{eq:kernel-floor-definition}.

\begin{theorem}[Floor for the rational four-plane]
\label{thm:kstar-floor}
For the rational four-plane $K^\ast$ reproduced in the supplement,
\[
  m(K^\ast)\geq \frac{1}{32}.
\]
\end{theorem}

\begin{proof}
Insert the four rational basis matrices into~\eqref{eq:spread-pencil}.  Exact
branch-and-bound, using the same rational interval and vertex-cone procedure
as above, proves that the normalized squared spread is at most $63/32$.
Equation~\eqref{eq:spread-pencil} then gives $m(K^\ast)\geq1/32$.  The finite
cover, consisting of $1\,786$ closed and $285$ pruned leaves, and the engine
recomputing its rational bounds are supplied in the supplement.
\end{proof}

The two designs may therefore be compared only at the level of their
proved lower-bound values.  The floor $1/32$ for $K^\ast$ is exactly
$32$ times the positive-projector floor $1/1024$.  Likewise,
\[
 \frac{1/(\pi\sqrt{32})}{\sqrt{L_0}/\pi}
 =\frac1{\sqrt{32L_0}}\approx64.3081.
\]
This is not a ratio of the unknown true minima or true stability constants.
The rational four-plane has the stronger proved guarantee, while the
positive-projector design retains the constructive tomography interpretation.

Exploratory floating-point optimization gives approximately
$m(\widehat K_{\rm proj})=0.0010212$ and $m(K^\ast)=0.10407$.
These are uncertified estimates, not proofs of global optimality. They
suggest that the projector floor is comparatively tight while the rational
floor is conservative; none of the inequalities above depends on them.

\begin{corollary}[An abstract isometric rotation-angle bound]
\label{cor:kstar-conditioning}
Let $E=K^{\ast\perp}\cap\mathcal S_0$ and
$H=\{x\in\RR^6:\sum_{j=1}^6x_j=0\}$.  Both spaces have dimension five,
so choose any linear isometry $T:E\to H$ and define
$\Psi_6(R(q))=T\Pi_E(qq^{\mathsf T}-I/4)$.
This is well defined under $q\sim-q$ and realizes pair differences by
\begin{equation}\label{eq:kstar-difference-realization}
 \Psi_6(R(q))-\Psi_6(R(r))
 =T\!\left(\Pi_E(qq^{\mathsf T}-rr^{\mathsf T})\right)
\end{equation}
for all unit quaternions $q,r$.  Define $L:\RR^6\to\RR^5$ by
\begin{equation}\label{eq:H-partial-isometry}
 L_{ij}=\begin{cases}
  1/\sqrt{i(i+1)},&j\leq i,\\
  -i/\sqrt{i(i+1)},&j=i+1,\\
  0,&j>i+1,
  \end{cases}
  \qquad 1\leq i\leq5,\quad 1\leq j\leq6.
\end{equation}
Then $LL^{\mathsf T}=I_5$ and $\|Lx\|_2=\|x\|_2$ for $x\in H$.  For the
five-coordinate map $\Phi^\ast=L\Psi_6$ and rotation-angle distance
$d_\angle$,
\[
 \|\Phi^\ast(R_1)-\Phi^\ast(R_2)\|
 \geq \frac{1}{\pi\sqrt{32}}d_\angle(R_1,R_2),
 \qquad
 c_{\Phi^\ast}\geq\frac{1}{\pi\sqrt{32}}>0.0562.
\]
The quaternion dictionary makes these fixed linear coordinates on $\SO(3)$;
the conclusion concerns this abstract $K^\ast$ construction, not a
positive-projector or device realization.
\end{corollary}

\begin{proof}
The rows of $L$ are orthonormal, and its kernel is the span of
$(1,1,1,1,1,1)$, proving the asserted isometry on $H$.  Hence
\eqref{eq:kstar-difference-realization} and
Theorem~\ref{thm:kstar-floor} make the realized objective equal to
$\|\Pi_EZ\|_F^2$ with lower bound $1/32$ on $\mathcal W$.  Apply
Lemma~\ref{lem:angle-reduction}.
\end{proof}

The distinction between injectivity and
stability in abstract phase retrieval, including compact-group formulations,
is treated in the primary literature
\cite{BandeiraCahillMixonNelson2014,XiaXuXu2024,AmirBendoryDymEdidin2025}.

\section{Fixed-linear optimality}\label{sec:optimality}

Wang and Xu define the minimal number $m_{\RR}(d)$ of generalized real
quadratic measurements required for phase retrieval and prove
$m_{\RR}(4)=6$ \cite[Theorem 5.1 and the following remark]{WangXu2019}.
The six rank-two projection construction used above originates in Xu's
low-rank-recovery and projection phase-retrieval construction
\cite{Xu2018}; Theorem~\ref{thm:kernel-certificate} supplies the exact
specialization needed here.
These generalized quadratic counts are not general tomography counts:
minimal qubit tomography and tomography under prior information use different
objects and measurement classes
\cite{RehacekEtAl2004,HeinosaariMazzarellaWolf2013}.

\begin{corollary}[Fixed-linear optimality]
\label{cor:fixed-five-optimality-external}
By the published theorem $m_{\RR}(4)=6$, five is the minimum
number of fixed real-linear functionals on the standard embedded
$\SO(3)\subset\RR^9$ that can be globally injective.
\end{corollary}

\begin{proof}
Suppose four fixed linear functionals were injective on $\SO(3)$.  Composing
them with the homogeneous quadratic map $q\mapsto R(q)$ gives four
generalized quadratic measurements on $\RR^4$.  Adjoin the norm measurement
$q^\mathsf Tq$.  If two nonzero vectors have the same five quadratic values,
their norms agree; after normalization, injectivity on $\SO(3)$ identifies
their unit quaternions up to sign.  Hence these five quadratics would perform
real phase retrieval in $\RR^4$, contradicting $m_{\RR}(4)=6$.
\end{proof}

The value $m_{\RR}(4)=6$ is a published input and is not reproved here.  The
conclusion concerns fixed real-linear functionals, not adaptive queries,
nonlinear encodings, primitive physical readings, or experimental cost.
There is also an independent topological obstruction: the compact space
$\SO(3)$, homeomorphic to $\mathbb{RP}^3$, does not embed in $\RR^4$, although it embeds
in $\RR^5$ \cite[Sec.~4.1]{ZhouEtAl2019}.  A continuous injective map
from a compact space into Euclidean space is an embedding, so this rules
out every continuous four-coordinate encoding, not only fixed linear
functionals.  The phase-retrieval reduction above gives a second proof
within the declared quadratic measurement model.

\section{Adaptive five-reading implementation comparison}
\label{sec:adaptive_comparison}

Here \emph{minimal} is relative to the column-first real-linear model below.
It does not transfer to POVM tomography, trace-one state estimation, or
arbitrary prior information
\cite{RehacekEtAl2004,HeinosaariMazzarellaWolf2013}.

\begin{definition}[Column-first adaptive scheme]\label{def:adaptive-protocol}
For $C=[u\ v\ z]\in\SO(3)$, first read all three entries of $u$.  Choose
$p\in\arg\max_i|u_i|$, taking the smallest index in a tie.  Then read the two components $v_j$ with
$j\ne p$, set
\begin{equation}\label{eq:adaptive-solve}
 v_p=-\frac{\sum_{j\ne p}u_jv_j}{u_p},
\end{equation}
 and complete $z=u\times v$.  The resulting matrix can be checked directly
 for unit columns, orthogonality, and positive orientation.
\end{definition}

\begin{proposition}[Scoped adaptive sufficiency and minimality]
\label{prop:adaptive-five}
The column-first scheme reconstructs every $C\in\SO(3)$ from five
support-aware readings.  After a full first column has been fixed, one
additional real-linear reading cannot be injective on the residual circle,
whereas the two selected nonpivot entries are injective.  Thus five is
necessary and sufficient within this declared column-first real-linear
scheme.
\end{proposition}

\begin{proof}
A unit vector $u$ has $|u_p|\geq1/\sqrt3$ at the selected
maximum-magnitude component.  Orthogonality
$u^\mathsf Tv=0$ gives~\eqref{eq:adaptive-solve}; the two nonpivot
coordinates therefore determine $v$ uniquely.  Exact checks certify
$u,v$ as an orthonormal pair, and $u\times v$ chooses the
orientation-preserving completion.

Once $u$ is fixed, write the residual column as
$v(\phi)=a\cos\phi+b\sin\phi$ in an orthonormal basis of $u^\perp$.  Any
one real-linear reading is $A\cos\phi+B\sin\phi$ and has an exact collision
on the circle.  By contrast, equality of the two nonpivot components and
orthogonality force equality of the pivot component, hence of $v$.
\end{proof}

\begin{proposition}[Conditional stability of the adaptive pivot]
\label{prop:adaptive-stability}
Suppose the first column $u$ is known exactly.  If the two measured
nonpivot coordinates of $v$ have error vector $e\in\mathbb R^2$, the
reconstruction~\eqref{eq:adaptive-solve}, before any normalization or
orthogonality correction, has error
\begin{equation}\label{eq:adaptive-stability}
 \|\delta v\|\leq \frac{\|e\|}{|u_p|}\leq\sqrt3\,\|e\|.
\end{equation}
Thus maximum-magnitude pivoting gives a uniform bound for this step
without another reading.
\end{proposition}
\begin{proof}
Since $\|u\|=1$, the selected pivot satisfies $|u_p|\geq1/\sqrt3$.
Writing $u_{-p}$ for its two nonpivot coordinates, the pivot error is
$\delta v_p=-u_{-p}^{\mathsf T}e/u_p$.  Cauchy--Schwarz gives
\[
 \|\delta v\|^2
 =\|e\|^2+\frac{|u_{-p}^{\mathsf T}e|^2}{|u_p|^2}
 \leq \left(1+\frac{1-|u_p|^2}{|u_p|^2}\right)\|e\|^2.
\]
\end{proof}
This is not a full noisy-instrument theorem: errors in $u$, pivot selection
from noisy first-column data, and subsequent completion are not covered.
Choosing the first nonzero component would preserve exact reconstruction
but allow an arbitrarily small denominator.  The supplement supplies an
exact-rational regression of the pivot rule and this distinction; the
general bound follows from the proof above, not from those examples.

\begin{table}[ht]
\centering
\small
\begin{tabular}{@{}p{0.28\textwidth}p{0.12\textwidth}p{0.49\textwidth}@{}}
\toprule
Instrument & Count & Mathematical scope\\
\midrule
Antisymmetric $w(C)$ & $3$ & Local chart on the positive-cosine branch\\
$w(C)$ plus trace & $4$ & Fails globally on every half-turn\\
Column-first scheme & $5$ & Global, exact, adaptive; minimal only in its
declared model\\
Functionals $\ell_1,\ldots,\ell_5$ & $5$ & Global, exact, fixed and
nonadaptive; quantitative positive-projector bounds; fixed-linear optimality
by $m_{\RR}(4)=6$\\
Rational $K^\ast$ coordinates & $5$ & Separate abstract design;
proved floor $1/32$ and abstract lower constant; not a projector or
physical implementation\\
\bottomrule
\end{tabular}
\caption{The adaptive result is an implementation comparison, not an open
substitute for the fixed theorem.}
\label{tab:tomography-comparison}
\end{table}

\section{Limitations and related directions}\label{sec:boundaries}

The reconstruction theorems concern exact finite-dimensional data;
Proposition~\ref{prop:adaptive-stability} adds only the stated conditional
perturbation bound.  The loops are prescribed
mathematical paths in $\Pos$, and their return action is relational in the
sense of~\eqref{eq:cotransported-readout}; no physical realization of an
arbitrary loop is asserted.  Likewise, the five fixed coordinates
are synthesized linear functionals, not five primitive experimental shots.
The fixed-linear lower bound uses the published identity
$m_{\RR}(4)=6$, and the adaptive minimality statement applies only to the
column-first model of Definition~\ref{def:adaptive-protocol}.

The conditioning estimates concern the declared exact coordinate maps.  They
do not establish stochastic-noise robustness, sample complexity, device
performance, or optimality among all four-planes.  In particular, the lower
angle constant for $K^\ast$ applies to the abstract realization
\eqref{eq:kstar-difference-realization}, not to an unspecified device.
Robust synchronization under
corruption is a natural comparison direction
\cite{Singer2011,BandeiraSingerSpielman2013,LangeEtAl2015,ElKarouiWu2016},
but incorporating noise into the Green-return reconstruction problem remains
open.

\section{Conclusion}\label{sec:conclusion}

The Green pairing determines the affine-invariant geometry of positive
Hermitian $2\times2$ instruments.  Its determinant-one shape leaf is
hyperbolic with the exact normalization $K=-1/2$, and its exact transport
turns closed instrument loops into relational $\SO(3)$ return actions.
Three readings are locally sufficient, the natural fourth trace reading does
not cross the half-turn wall, and five fixed synthesized functionals are
globally injective.  Their proof is constructive:
six exact rank-two projections, exact whitening, an exact
quaternion--rotation dictionary, and a complete projective real-kernel
exclusion.  Fixed-linear optimality remains explicitly dependent on the
published value $m_{\RR}(4)=6$.  The adaptive five-reading scheme remains
an elementary alternative within its narrower information model.

The positive-projector map has
$m(\widehat K_{\rm proj})\geq1/1024$, $\min F_{\rm proj}\geq L_0$, the
two-sided estimate~\eqref{eq:projector-angle-bounds}, and
$\kappa_{\rm proj}\leq560.6809131191\ldots$.  The separate rational plane
$K^\ast$ has $m(K^\ast)\geq1/32$ and, for the constructed isometric
coordinates, a lower angle constant exceeding $0.0562$.  The factor $32$ in
the floor bounds and the factor $64.3081\ldots$ in the corresponding proved
lower constants compare guarantees, not unknown true optima.  Determining the
best four-plane, and understanding which well-conditioned abstract designs
admit natural positive or physical realizations, are the principal open
problems suggested by this analysis.

\section*{Acknowledgment of assistance}
AI-assisted research, computation, checking, and editorial support used
OpenAI Codex and Anthropic Claude. Luqman Mahmood directed the work
and is responsible for the claims and final text. These internal
checks are not external peer review.

\begingroup
\small
\bibliographystyle{plain}
\bibliography{references}
\endgroup

\end{document}
